% Bagian: first-order-logic
% Bab: syntax-and-semantics
% Subbagian: free-vars-sentences

\documentclass[../../../include/open-logic-section]{subfiles}

\begin{document}

\olfileid[id]{fol}{syn}{fvs}

\olsection{\printtoken{P}{variable} Bebas dan \printtoken{P}{sentence}}

\begin{defn}[Kemunculan bebas !!a{variable}]
\ollabel{defn:free-occ}
Kemunculan \emph{bebas} !!a{variable} dalam !!a{formula} didefinisikan
secara induktif sebagai berikut:
\begin{enumerate}
\item \indcase*{!A}{$!A$ atomik}{semua kemunculan !!{variable} dalam
  $\indfrm$ adalah bebas.}

\tagitem{prvNot}{\indcase{!A}{\lnot !B}{kemunculan bebas !!{variable}
  dalam $\indfrm$ tepat sama dengan kemunculan bebas dalam $!B$.}}{}

\item \indcase{!A}{(!B \ast !C)}{kemunculan bebas !!{variable} dalam
  $\indfrm$ ialah kemunculan bebas dalam $!B$ bersama kemunculan bebas
  dalam~$!C$.}

\tagitem{prvAll}{\indcase{!A}{\lforall[x][!B]}{kemunculan bebas
  !!{variable} dalam $\indfrm$ ialah semua kemunculan bebas dalam~$!B$
  kecuali kemunculan~$x$.}}{}

\tagitem{prvEx}{\indcase{!A}{\lexists[x][!B]}{kemunculan bebas
  !!{variable} dalam $\indfrm$ ialah semua kemunculan bebas dalam~$!B$
  kecuali kemunculan~$x$.}}{}
\end{enumerate}
\end{defn}

\begin{defn}[Variabel Terikat]
Suatu kemunculan !!a{variable} dalam formula~$!A$ disebut \emph{terikat}
jika kemunculan itu tidak bebas.
\end{defn}

\begin{prob}
Berikan definisi induktif bagi kemunculan variabel terikat dengan mengikuti
pola \olref[fol][syn][fvs]{defn:free-occ}.
\end{prob}

\begin{defn}[Lingkup]
\iftag{prvAll}{Jika $\lforall[x][!B]$ merupakan suatu kemunculan subformula
  dalam formula~$!A$, kemunculan $!B$ yang bersesuaian dalam~$!A$ disebut
  \emph{lingkup} dari kemunculan $\lforall[x]$ yang bersesuaian.
  \iftag{prvEx}{Demikian pula untuk $\lexists[x]$.}{}}{Jika
  $\lexists[x][!B]$ merupakan suatu kemunculan subformula dalam
  formula~$!A$, kemunculan $!B$ yang bersesuaian dalam~$!A$ disebut
  \emph{lingkup} dari kemunculan $\lexists[x]$ yang bersesuaian.}

Jika $!B$ merupakan lingkup dari suatu kemunculan kuantor
\iftag{prvAll}{$\lforall[x]$\iftag{prvEx}{ atau
    $\lexists[x]$}{}}{$\lexists[x]$} dalam~$!A$, maka kemunculan bebas $x$
dalam~$!B$ menjadi terikat dalam
\iftag{prvAll}{$\lforall[x][!B]$\iftag{prvEx}{ dan
$\lexists[x][!B]$}{}}{$\lexists[x][!B]$}. Kita mengatakan bahwa
kemunculan-kemunculan tersebut \emph{diikat oleh} kemunculan kuantor yang
disebutkan.
\end{defn}

\begin{ex}
Perhatikan formula berikut:
\[
\lexists[\Obj v_0][\underbrace{\Atom{\Obj A^2_0}{\Obj v_0,\Obj v_1}}_{!B}]
\]
$!B$ menyatakan lingkup $\lexists[\Obj v_0]$.
Kuantor tersebut mengikat kemunculan $\Obj v_0$ dalam $!B$, tetapi tidak
mengikat kemunculan $\Obj v_1$. Jadi, dalam kasus ini $\Obj v_1$ merupakan
variabel bebas.


Sekarang kita dapat melihat cara kerjanya dalam suatu !!{formula}~$!A$
yang lebih rumit:
\[
\lforall[\Obj v_0][\underbrace{(\Atom{\Obj A^1_0}{\Obj v_0} \lif
    \Atom{\Obj A^2_0}{\Obj v_0, \Obj v_1})}_{!B}] \lif \lexists[\Obj
  v_1][\underbrace{(\Atom{\Obj A^2_1}{\Obj v_0, \Obj v_1} \lor \lforall[\Obj v_0][\overbrace{\lnot \Atom{\Obj A^1_1}{\Obj v_0}}^{!D}])}_{!C}]
\]
$!B$ merupakan lingkup dari $\lforall[\Obj v_0]$ pertama, $!C$ merupakan
lingkup dari $\lexists[\Obj v_1]$, dan $!D$ merupakan lingkup dari
$\lforall[\Obj v_0]$ kedua. $\lforall[\Obj v_0]$ pertama mengikat
kemunculan-kemunculan $\Obj v_0$ dalam~$!B$, $\lexists[\Obj v_1]$ mengikat
kemunculan $\Obj v_1$ dalam $!C$, dan $\lforall[\Obj v_0]$ kedua mengikat
kemunculan $\Obj v_0$ dalam~$!D$. Kemunculan pertama $\Obj v_1$ dan
kemunculan keempat $\Obj v_0$ bebas dalam~$!A$. Kemunculan terakhir
$\Obj v_0$ bebas dalam $!D$, tetapi terikat dalam $!C$ dan~$!A$.
\end{ex}

\begin{defn}[Kalimat]
!!^a{formula}~$!A$ merupakan \article{sentence} \emph{!!{sentence}} jika dan
hanya jika formula itu tidak memuat kemunculan bebas !!{variable}s.
\end{defn}

% Tambahkan contoh.

\end{document}

